Una matriu quadrada $A$ és \emph{nilpotent} si alguna de les seves potències és la matriu nul·la. Si $n$
és el menor enter positiu tal que $A^n=0$, es diu que $A$ és nilpotent de grau $n$.

\begin{apartats}

\apartat[1,25]{0,75}
Demostra que la matriu $N=\begin{pmatrix}0&1&0\\-1&0&1\\0&1&0\end{pmatrix}$ és nilpotent de grau 3.

\begin{solucio}
\[
N^2=\begin{pmatrix}-1&0&1\\0&0&0\\-1&0&1\end{pmatrix}\ne0,\qquad
N^3=N^2N=\begin{pmatrix}0&0&0\\0&0&0\\0&0&0\end{pmatrix}.
\]
$N^2\ne0$ i $N^3=0$: $N$ és nilpotent de grau 3.
\end{solucio}

\begin{tria}{nilpotents}
\itemtria{grau-2}{1}{1,25}
Troba totes les matrius de la forma $B=\begin{pmatrix}a&1\\b&-a\end{pmatrix}$ que són nilpotents de
grau 2.

\begin{solucio}
\[
B^2=\begin{pmatrix}a&1\\b&-a\end{pmatrix}\begin{pmatrix}a&1\\b&-a\end{pmatrix}
=\begin{pmatrix}a^2+b&0\\0&a^2+b\end{pmatrix}.
\]
$B^2=0$ si i només si $b=-a^2$. Com que $B\ne0$ (té un 1), són nilpotents de grau 2 totes les matrius
$B=\begin{pmatrix}a&1\\-a^2&-a\end{pmatrix}$, amb $a\in\mathbb R$.
\end{solucio}

\itemtria{sense-calcular}{1}{1,25}
Sense fer cap producte més, digues quant valen $N^4$ i $N^{100}$. Després, calcula $(I+N)^3$ fent servir
que $N^3=0$.

\begin{solucio}
$N^4=N^3N=0$, i per a qualsevol $k\ge3$, $N^k=N^3N^{k-3}=0$: també $N^{100}=0$.\\
Com que $I$ commuta amb $N$, el cub d'una suma es desenvolupa com amb nombres:
$(I+N)^3=I+3N+3N^2+N^3=I+3N+3N^2$.
\[
(I+N)^3=\begin{pmatrix}1&0&0\\0&1&0\\0&0&1\end{pmatrix}+\begin{pmatrix}0&3&0\\-3&0&3\\0&3&0\end{pmatrix}
+\begin{pmatrix}-3&0&3\\0&0&0\\-3&0&3\end{pmatrix}=\begin{pmatrix}-2&3&3\\-3&1&3\\-3&3&4\end{pmatrix}.
\]
\end{solucio}
\end{tria}

\begin{nomesllarg}
\apartat{0,75}
Comprova que $(I-N)(I+N+N^2)=I$.

\begin{solucio}
$(I-N)(I+N+N^2)=I+N+N^2-N-N^2-N^3=I-N^3=I$, perquè $N^3=0$.
\end{solucio}
\end{nomesllarg}

\end{apartats}
